% Canonical Paper II source.
Throughout this section
\[
 \mathfrak E=(M,g,a;\mu,\lambda)
\]
is a regular arithmetic expression space in the sense of Paper~I, with real
constants $\mu,\lambda$.  Thus $M$
is an oriented smooth surface, possibly with boundary, $g$ is a smooth
Riemannian metric, $a\in C^\infty(M,\R)$, $\mu\ne0$, and
\begin{equation}\label{eq:paper2-regular-eikonal}
 |\nabla a|_g^2=\mu^2+\lambda^2a^2.
\end{equation}
The metric and orientation are part of the imported structure.  In particular,
they are not reconstructed from an analytic operator introduced later.

\subsection{The arithmetic frame and its structure functions}

Write
\[
 q=(\mu^2+\lambda^2a^2)^{1/2}.
\]
Since $\mu\ne0$, $q$ never vanishes.  Paper~I therefore supplies the unique
positively oriented orthonormal frame $(\Xu,\Xv)$ satisfying
\begin{equation}\label{eq:paper2-assignment-frame-rules}
 \Xu a=\mu,
 \qquad
 \Xv a=\lambda a.
\end{equation}
No coordinate interpretation is implicit here: the frame is generally
noncommuting.  Define its smooth structure functions
$c_u,c_v\in C^\infty(M)$
by
\begin{equation}\label{eq:surface-frame-bracket}
 [\Xu,\Xv]=c_u\Xu+c_v\Xv.
\end{equation}
Applying this identity to the assignment gives the constraint
\begin{equation}\label{eq:surface-structure-assignment-constraint}
 \mu c_u+\lambda a c_v=\mu\lambda.
\end{equation}
Equation~\eqref{eq:surface-structure-assignment-constraint} is only one scalar
relation.  It does not determine $c_u,c_v$, the Gaussian curvature, or
$\Delta_ga$ on a general regular AES.

Let $d\vol_g$ be the positive Riemannian measure.  The structure functions
also determine the divergences of the arithmetic frame.

\begin{proposition}[Structure functions and divergence]
\label{prop:surface-frame-divergence}
The arithmetic frame satisfies
\begin{equation}\label{eq:surface-frame-divergences}
 \diver_g\Xu=-c_v,
 \qquad
 \diver_g\Xv=c_u.
\end{equation}
Consequently, on compactly supported smooth scalar fields in the interior
$M^\circ=M\setminus\partial M$, its formal adjoints are
\begin{equation}\label{eq:surface-frame-formal-adjoints}
 \Xu^*=-\Xu+c_v,
 \qquad
 \Xv^*=-\Xv-c_u.
\end{equation}
Here the star denotes a formal adjoint on
$C_c^\infty(M^\circ)\subset L^2(M,d\vol_g)$, not a choice of a closed
Hilbert-space realization.
\end{proposition}

\begin{proof}
For a positively oriented orthonormal frame on a surface, metric compatibility
and vanishing torsion give
\[
 \langle\nabla_{\Xv}\Xu,\Xv\rangle=-c_v,
 \qquad
 \langle\nabla_{\Xu}\Xv,\Xu\rangle=c_u.
\]
The remaining diagonal terms vanish because the frame has unit length.  Taking
the trace of $Y\mapsto\nabla_Y\Xu$ and of
$Y\mapsto\nabla_Y\Xv$ proves
\eqref{eq:surface-frame-divergences}.  If $X$ is a real vector field and
$f,h\in C_c^\infty(M^\circ;\C)$, the divergence theorem gives
\[
 \int_M (Xf)\,\overline h\,d\vol_g
 =-\int_M f\,\overline{(X+\diver_gX)h}\,d\vol_g.
\]
Substituting~\eqref{eq:surface-frame-divergences} yields
\eqref{eq:surface-frame-formal-adjoints}.  A coordinate-free sign check is
recorded in Appendix~\ref{app:frame-calculations}.
\end{proof}

\subsection{The Laplace--Beltrami expression and energy}

We retain the divergence-of-gradient convention
\begin{equation}\label{eq:surface-laplacian-convention}
 \Delta_gf=\diver_g(\nabla f).
\end{equation}
With this convention $\Delta_g$ is nonpositive in the $L^2$ quadratic-form
sense on compactly supported functions.

\begin{proposition}[Arithmetic-frame Laplacian]
\label{prop:surface-frame-laplacian}
For every $f\in C^\infty(M)$, the following pointwise identity holds:
\begin{equation}\label{eq:surface-frame-laplacian}
 \boxed{
 \Delta_gf
 =\Xu^2f+\Xv^2f-c_v\Xu f+c_u\Xv f.
 }
\end{equation}
On $C_c^\infty(M^\circ)$, the same differential expression has the formal
energy factorization
\begin{equation}\label{eq:surface-real-energy-factorization}
 -\Delta_g=\Xu^*\Xu+\Xv^*\Xv.
\end{equation}
\end{proposition}

\begin{proof}
Because the frame is orthonormal,
\[
 \nabla f=(\Xu f)\Xu+(\Xv f)\Xv.
\]
Using
$\diver(\phi X)=X\phi+\phi\diver X$ and
Proposition~\ref{prop:surface-frame-divergence} gives
\eqref{eq:surface-frame-laplacian}.  Substitution of the formal adjoints
\eqref{eq:surface-frame-formal-adjoints} gives
\eqref{eq:surface-real-energy-factorization} directly.
\end{proof}

In particular, the raw sum of squares
\begin{equation}\label{eq:surface-raw-sum-squares}
 S_M:=\Xu^2+\Xv^2
\end{equation}
is not the Laplace--Beltrami operator unless
$-c_v\Xu+c_u\Xv=0$.  The drift in
\eqref{eq:surface-frame-laplacian} is fixed by $d\vol_g$; it is not an
optional lower-order perturbation.

For later operator statements, define the densely defined nonnegative form
\begin{equation}\label{eq:surface-dirichlet-preform}
 \mathcal E_0(f,h)
 :=\int_M
 \bigl((\Xu f)\overline{\Xu h}
       +(\Xv f)\overline{\Xv h}\bigr)\,d\vol_g,
 \qquad
 f,h\in C_c^\infty(M^\circ;\C).
\end{equation}
Its form norm is
\[
 \|f\|_{\mathcal E}^2
 :=\|f\|_{L^2(d\vol_g)}^2+\mathcal E_0(f,f).
\]

\begin{definition}[Minimal energy domain and Dirichlet realization]
\label{def:surface-energy-domain}
Let $H_0^1(M,g)$ denote the completion of
$C_c^\infty(M^\circ)$ in the form norm above.  The closure
$\mathcal E$ of $\mathcal E_0$ is the Dirichlet energy form.  We denote by
$-\Delta_{g,D}$ the nonnegative self-adjoint operator associated with
$\mathcal E$ by the representation theorem for closed forms.
\end{definition}

This notation makes no assertion that the minimal differential operator is
essentially self-adjoint.  If $M$ has boundary, it selects the Dirichlet form
domain; Neumann, Robin, and mixed realizations require different form domains.
If $M$ is complete without boundary, standard completeness results may identify
this realization with the unique self-adjoint closure, but no such conclusion is
needed for the local identities below.

\begin{proposition}[Weak and operator meanings]
\label{prop:surface-weak-operator-meanings}
For $f\in C_c^\infty(M^\circ)$,
\begin{equation}\label{eq:surface-green-identity-compact-support}
 \mathcal E_0(f,h)
 =\langle-\Delta_gf,h\rangle_{L^2(d\vol_g)}
 \qquad
 (h\in C_c^\infty(M^\circ)).
\end{equation}
More generally, $f\in\Dom(-\Delta_{g,D})$ precisely when there is
$k\in L^2(M,d\vol_g)$ such that
\[
 \mathcal E(f,h)=\langle k,h\rangle_{L^2(d\vol_g)}
 \quad\text{for every }h\in H_0^1(M,g),
\]
in which case $-\Delta_{g,D}f=k$.
\end{proposition}

\begin{proof}
Equation~\eqref{eq:surface-green-identity-compact-support} follows from
\eqref{eq:surface-real-energy-factorization} and integration by parts; compact
support removes all boundary terms.  The second statement is the defining
characterization of the operator represented by the closed form
$\mathcal E$.
\end{proof}

Thus three levels will be kept separate throughout the paper:
\begin{enumerate}
 \item pointwise identities of differential expressions on $C^\infty$;
 \item formal-adjoint identities on $C_c^\infty(M^\circ)$;
 \item statements about a closed operator on a declared Hilbert-space domain.
\end{enumerate}
Only the third level supports spectral terminology.
